\section*{Coda: information constraints survive architecture}
\addcontentsline{toc}{section}{Coda: information constraints survive architecture}

Capacity and admissible information impose different limits.  A router can predict an edge
accurately by reading its endpoints, while strict conditional semantics require its selection to
survive changes to those endpoints.  The Bayes floor quantifies the resulting inference limit.
A teacher can improve training; only additional inference-time evidence or a weaker invariance
contract can lower that floor.

The two realizability chapters therefore give separate tests for representability, execution,
and measurability.  A useful dynamic route can fail the strict conditional test without failing
its computational task.  The final empirical question is which declared operator corresponds to
evolutionary constraint, structure, or a measured assay response.  The next chapter specifies
experiments that can reject those correspondences individually.
